\documentclass[11pt,letterpaper]{article}
\usepackage[margin=0.84in]{geometry}
\usepackage[T1]{fontenc}
\usepackage{lmodern}
\usepackage{amsmath,amssymb,amsthm,mathtools,booktabs,array}
\usepackage{xcolor,microtype,xurl,hyperref,fancyhdr}
\hypersetup{colorlinks=true,linkcolor=blue!50!black,citecolor=blue!50!black,urlcolor=blue!50!black,
pdftitle={Recovering signs from a binomial digit},
pdfauthor={AI-assisted research for Vladimir Reshetnikov and ProveIt},
pdfsubject={An explicit quartic consistency equation in 247 integer variables}}
\newtheorem{theorem}{Theorem}
\newtheorem{lemma}[theorem]{Lemma}
\newtheorem{proposition}[theorem]{Proposition}
\newtheorem{remark}[theorem]{Remark}
\newcommand{\ZZ}{\mathbb Z}
\newcommand{\NN}{\mathbb N}
\newcommand{\Con}{\operatorname{Con}}
\newcommand{\ZFC}{\mathsf{ZFC}}
\newcommand{\ZF}{\mathsf{ZF}}
\newcommand{\Ccal}{\mathsf{C222}}
\newcommand{\G}{\mathcal G}
\newcommand{\file}[1]{\texttt{\nolinkurl{#1}}}
\newcommand{\FinalTerms}{4,755}
\newcommand{\FinalHeight}{8,815}
\newcommand{\FinalRows}{182}

\pagestyle{fancy}
\fancyhf{}
\lhead{\small Recovering signs from a binomial digit}
\rhead{\small ProveIt research contribution}
\cfoot{\thepage}
\setlength{\headheight}{14pt}
\setlength{\emergencystretch}{2em}
\allowdisplaybreaks[2]
\title{\textbf{Recovering signs from a binomial digit}\\[0.5em]
\large An explicit quartic consistency equation in 247 integer variables}
\author{Research prepared with AI assistance\\
\small For Vladimir Reshetnikov and the ProveIt repository}
\date{25 October 2026}
\begin{document}
\maketitle
\begin{abstract}
We give a completely specified ordinary polynomial
$P\in\ZZ[z_1,\ldots,z_{247}]$ of degree four, with \FinalTerms{} collected
monomials and coefficient height \FinalHeight{}, such that ZFC proves
that $P$ has no integer zero if and only if ZFC is consistent. The polynomial
is the sum of squares of \FinalRows{} explicitly printed quadratic
residuals. Its mathematical reduction is proved by composition with the
retained, pinned C222 refutation calibration.

The new step recovers the sign of the factorial binomial digit before
using the usual ceiling-sign certificates. A retained denominator-ratio
equation forces an early exponent bound, and the self-sized radix then
separates the positive binomial digit from all signed alternatives.
Three scalar nonnegativity certificates can consequently be replaced by
one product certificate and a cofinally available square restriction on
the radix. Combined with an inherited deterministic Pell comparison
center, this removes five variables from the 252-integer common-remainder
predecessor. Certified integral changes of coordinates and residual basis
give the final height and support.

This is a measured improvement of a particular small-coefficient direct
quartic, not a variable-count record or a proof of global minimality.
The new unbounded argument is an ordinary mathematical proof; a fresh
end-to-end proof-assistant formalization is not claimed. Independent
finite checks validate the exact source changes, integral restorations,
all coefficients, and every residual printed here.
\end{abstract}
\clearpage
\tableofcontents
\clearpage

\section{The delivered equation and the meaning of simplicity}
\label{sec:equation}

The deliverable is the single equation
\begin{equation}\label{eq:main}
 \boxed{\quad P(z_1,\ldots,z_{247})
       :=\sum_{j=1}^{182}R_j(z_1,\ldots,z_{247})^2=0,
       \qquad (z_1,\ldots,z_{247})\in\ZZ^{247}.\quad}
\end{equation}
Appendix~\ref{app:rows} prints every $R_j$. The file
\file{equations/quartic247/equation.json} gives the same finite list, and
\file{equations/quartic247/expanded_polynomial.json} gives every collected
coefficient of $P$. All variables are unrestricted integers. There are
no suppressed parameters, variable-exponent operations, factorial
operations, proof predicates, or quantified helper functions in
\eqref{eq:main}. The notation for such objects below is used only to
prove the meaning of ordinary arithmetic registers.

\begin{theorem}\label{thm:main}
Let $P$ be the literal polynomial \eqref{eq:main}. With the syntactic
axiomatization of ZFC used by the retained ProveIt calibration, ZFC proves
\[
 \bigl(\exists z_1,\ldots,z_{247}\in\ZZ,\ P(z_1,\ldots,z_{247})=0\bigr)
 \ \longleftrightarrow\
 \text{there is a finite ZFC refutation}.
\]
Equivalently,
\[
 \ZFC\vdash
 \left[
 \bigl(\forall \mathbf z\in\ZZ^{247},\ P(\mathbf z)\ne0\bigr)
 \longleftrightarrow \Con(\ZFC)
 \right].
\]
Its degree is exactly $4$, its support has exactly $4\,755$ monomials,
and its maximum absolute coefficient is $8\,815$.
\end{theorem}

Here consistency means the absence of a finite syntactic derivation of
contradiction in the fixed first-order ZFC theory. We do not assert
either side of the equivalence separately. The theorem does not
conflict with the second incompleteness theorem: if ZFC is consistent,
the construction does not supply a ZFC proof of the universal
nonvanishing statement.

\begin{center}
\small
\begin{tabular}{@{}lrrrrr@{}}
\toprule
Endpoint over unrestricted integers & Variables & Rows & Degree & Terms & Height\\
\midrule
Common-remainder predecessor &252&182&4&5,340&128,446\\
New source after integral projection &247&182&4&5,310&128,446\\
Delivered coordinate and residual basis &247&182&4&4,755&8,815\\
Prior shared-residue comparison &260&186&4&4,770&9,191\\
\bottomrule
\end{tabular}
\end{center}

The comparison is pinned to repository commit
\file{3ecd6a6cf53d47519228328567f45fb7ee9781d9}, including the incoming
archives listed in \file{SOURCE_LEDGER.json}. Relative to the matched
252-variable predecessor, the delivered polynomial removes five
variables and 585 collected monomials, and lowers the height by
119,631. Relative to the prior 260-variable small-height endpoint, it
also improves all three displayed size measures.

There is no canonical total ordering of Diophantine equations by
simplicity. The repository already contains a quartic in 37 unrestricted
integers with large fixed coefficients, as well as still smaller
variable counts at larger degrees \cite{repo}. Those are not displaced
as variable-count achievements by this result. Our concrete objective
is a direct, fully expanded, modest-coefficient quartic with transparent
arithmetic witnesses. The route uses the numerical calculus C222
directly; it does not encode a K215 computation or instantiate Jones's
universal polynomial. C222's retained logical calibration is nevertheless
an essential inherited dependency.

For small integrity checks on the literal expansion, the constant
coefficient is $3494$, the coefficient of $z_1^4$ is $13$, and the
coefficient of $z_{12}^2z_{14}^2$ is $8815$. These are not a substitute
for the complete coefficient file.

\section{The retained arithmetic interface}
\label{sec:interface}

We begin with the direct common-remainder construction \cite{common}.
Its semantic registers satisfy the following ordinary identities before
their intended signs have been established:
\begin{align}
 t&=16+h_0^2,                  & B&=r+s+1,                    \label{eq:fronta}\\
 BF&=V+s,                     & y&=V(V-1),                   \label{eq:frontb}\\
 q&=B+o+1,                    & x&=1+[1+(t+1)y]u_t,           \label{eq:frontc}\\
 B\mu&=F^2+C_0+1,             & H_0&=1+\sum_{i=1}^{9}a_i.    \label{eq:frontd}
\end{align}
The $a_i$ are the nine actual encoded representatives, including $C_0$.
They are not initially presumed nonnegative. Composite expressions in
these displays are represented by the existing quadratic circuit.

For integers $\alpha,\beta,\gamma$, put
\[
 \G(\alpha,\beta,\gamma)=\alpha^2+\beta^2+\gamma(\gamma+1).
\]
This is always nonnegative and represents every nonnegative integer.
Indeed, the three-square theorem applies to $4n+1$; in a representation
by three squares exactly one coordinate is odd. Writing it as
$2\gamma+1$ and the other two as $2\alpha,2\beta$ gives the assertion.
The predecessor uses three separate such representations for $r,s,o$.

The new source deletes these nine private coordinates and imposes only
\begin{equation}\label{eq:new}
 \boxed{\quad
 rs=\G(\alpha,\beta,\gamma),\qquad q=(6\zeta+3)^2.
 \quad}
\end{equation}
There are four fresh unrestricted integer coordinates here. We use
the weak product condition $rs\ge0$; an unnecessary strict condition
would change the polynomial. Both equations are quadratic in the
actual source registers.

\subsection{Power meanings are established before the missing signs}

The retained source has eighteen genuine Pell power requests, including
\[
 V=(F^2+1)^t,\quad N=B^x,\quad q=(BH_0)^x,\quad
 D_t=q^t,\quad D_{t+1}=q^{t+1},\quad D_y=q^{y+1},\quad M=q^x,
\]
the factorial numerator $U_f=(q+1)^{F^2+t}$, and the representative
numerator requests. These displays describe the meanings to be proved,
not extra operations in the polynomial.

The signed frame, the squared semantic cap, and the canonical
orientation argument are inherited from the ordinary proofs
\file{PROOF_315.md}, \file{ORIENTATION_FIRST_PROOF.md}, and
\file{STRICT_BUDGET_PROOF.md}, retained under
\file{source/inherited_common/upstream/round17/direct/}.
Before any arithmetic sign is used, they give
\[
 W\ge258,\qquad |\text{every power atom}|<W,\qquad
 |A|>W^W,\qquad |A-1|>4W^3.
\]
They bound the absolute classified Pell indices, including at negative
centers and negative nominal indices.

We retain all fourteen orientation terms. For a principal pair
$(X_i,Y_i)$ with requested index $k_i$, the pool has the form
\begin{equation}\label{eq:pool}
 \sum_{i=1}^{14}\bigl((Y_i-k_i)^2+1-X_i\bigr)=(A-1)Q.
\end{equation}
The actual summands $1-X_i$ need not be nonnegative. The proof reduces
them modulo $A-1$: the small abscissa residues are
$s_i\in\{-1,1\}$, and the small ordinate residues are signed classified
indices $b_i$. Thus the integer
\[
 E=\sum_{i=1}^{14}(b_i-k_i)^2+\sum_{i=1}^{14}(1-s_i)
\]
is nonnegative, divisible by $A-1$, and satisfies
$E<4W^3<|A-1|$. It is zero. This recovers the full canonical pairs.
The smaller selective nine-term pool from a previous packet is
\emph{not} substituted: its later sign bootstrap would not justify
the new order of argument.

\subsection{The inherited deterministic comparison center}

We combine \eqref{eq:new} with the already proved deterministic-center
construction \cite{center}. The common Pell pair satisfies
\[
 H^2-(A^2-1)J^2=1,\qquad J=C^2(2e+1).
\]
Replacing the old free comparison-center constructor by
$f=(1-A)H$ gives
\[
 D=A+Hf
  =A+(1-A)H^2
  =1-(A-1)(A^2-1)J^2.
\]
Consequently $D\equiv A\pmod H$ and $D\equiv1\pmod C$. The comparison
ordinate congruence can be $R=h+Cq_c$, without the old factor $4$.
This deletes one additional private coordinate. It is valid for signed
$D$, which is often negative.

The exact proof, retained as \file{proofs/PELL_CENTER_PRIOR.md}, uses
squared-ordinate divisibility and abscissa residue separation.
At an actual anchor index $i\ge2$, Pell growth gives
$i+|h|<|C|$; the congruence $i\equiv\pm h\pmod{|C|}$ therefore yields
$i=|h|$. At the exceptional unit ordinate $|C|=1$, the squared cap
forces $h=0$ and all other actual indices to vanish; the anchor then
contributes a positive term to the bounded orientation residue $E$,
a contradiction. The proof works with the full fourteen-term pool
used here. Completeness chooses an even semantic anchor and rebuilds
the private comparison witnesses.

\subsection{The normalized digit equations}

The retained common-remainder block contains the literal ratio rows
\begin{equation}\label{eq:ratios}
 D_t v_t=M,\qquad D_{t+1}v_{t+1}=M,\qquad D_yv_y=M,
\end{equation}
and the factorial numerator row
\begin{equation}\label{eq:digit}
 v_t U_f=M(B+qk_f)+R_f.
\end{equation}
The product $qk_f$ is already a circuit register. All eleven
normalized remainders satisfy the one norm identity
\begin{equation}\label{eq:norm}
 M^2=1+4\sum_{i=1}^{11}R_i^2+\G(\alpha_*,\beta_*,\gamma_*).
\end{equation}
The starred coordinates are independent of the new coordinates in
\eqref{eq:new}. Equation \eqref{eq:norm} implies $|R_i|<M/2$ once
$M>0$. The retained normalized-remainder proof supplies stronger
$1/32$ margins for completeness.

\section{The new step: recover the digit, then its signs}
\label{sec:new}

\begin{lemma}[Canonical powers at signed bases]\label{lem:signed}
Once the retained full orientation theorem has been applied, a power
request at a base $n$ with $|n|\ge2$ forces its requested index $k$
to be nonnegative and its output to be $n^k$.
\end{lemma}
\begin{proof}
The ordinary power graph uses the modulus
$T=2An-n^2-1$. The frame gives
$|T|>W^W+W$, and $T\equiv-1\pmod n$, so $n$ is invertible modulo
$T$. For $k\ge0$, the canonical Pell recurrence yields
$m\equiv n^k\pmod T$. The bound on all atoms and indices makes the
absolute difference smaller than $|T|$, hence $m=n^k$.
For $k=-j<0$, the inverse recurrence instead gives
$mn^j\equiv1\pmod T$. The same bound would imply $mn^j=1$,
impossible for $|n|\ge2$ and $j\ge1$.
The argument uses absolute values throughout, so negative $n$ causes
no change. At $n=1$ and a known nonnegative index the output is $1$;
at $n=-1$ it is $(-1)^k$. We will exclude base zero before using the
self-radix request.
\end{proof}

\begin{lemma}[An early exponent bound]\label{lem:early}
Every integer zero of the new source satisfies
\[
 q\ge9,\quad M=q^x>0,\quad x\ge t+1\ge17,\quad
 V\ge1,\quad y\ge0,\quad a_i\ge0,\quad H_0\ge1,\quad B\ne0.
\]
\end{lemma}
\begin{proof}
The square in \eqref{eq:new} is the square of an integer congruent
to $3$ modulo $6$, hence $q\ge9$. The genuine $q^x$ request and
Lemma~\ref{lem:signed} imply $x\ge0$ and $M=q^x$.
Since $t\ge16$, the $t+1$ request is also exact. The middle row of
\eqref{eq:ratios} is therefore
\[
 q^{t+1}v_{t+1}=q^x.
\]
If $x<t+1$, cancellation would make a multiple of $q$ equal to $1$.
Thus $x\ge t+1$. This uses no factorial identity or ceiling sign.

The known nonnegative $t$ and the base $F^2+1\ge1$ give
$V=(F^2+1)^t\ge1$, hence $y=V(V-1)\ge0$.
Every actual representative is an exponent in a genuine request
at $q+1\ge10$, so each is nonnegative. In particular $C_0\ge0$
and $H_0\ge1$. Finally
$B\mu=F^2+C_0+1>0$ excludes $B=0$ without assigning a sign to $\mu$.
\end{proof}

\begin{lemma}[The self-radix bounds a signed output]\label{lem:radix}
At every integer zero,
\[
 q=(BH_0)^x=|BH_0|^x\ge2^x,\qquad
 B^2\le q,\qquad |B|<q/2.
\]
\end{lemma}
\begin{proof}
The base $BH_0$ is nonzero. If its absolute value were $1$, its
canonical output at the already nonnegative index $x$ would be
$1$ or $-1$, contradicting $q\ge9$. Thus its absolute value is at
least $2$ and Lemma~\ref{lem:signed} applies. Positivity of $q$
gives the equality with $|BH_0|^x$. Since $x\ge17$ and $H_0\ge1$,
we have $B^2\le q$. Finally $\sqrt q<q/2$ for $q\ge9$.
\end{proof}

\begin{lemma}[A uniform coefficient bound]\label{lem:coefficient}
Put $m=F^2$ and $n=m+t$. Then
\[
 n\le x-1,\qquad 2^n\le q/2.
\]
\end{lemma}
\begin{proof}
The coefficient $1+(t+1)y$ in \eqref{eq:frontc} is positive.
Since $x\ge17$, its integer quotient $u_t$ is at least $1$, giving
\begin{equation}\label{eq:xbudget}
 x\ge 2+(t+1)y.
\end{equation}
If $m=0$, then $n=t\le x-1$ by Lemma~\ref{lem:early}.
If $m\ge1$, Bernoulli's inequality yields
\[
 V=(m+1)^t\ge1+mt\ge m+t=n.
\]
Here $V\ge2$, so $y=V(V-1)\ge V\ge n$, and
\eqref{eq:xbudget} is stronger than the required bound.
Now $q\ge2^x$ proves $2^n\le2^{x-1}\le q/2$.
\end{proof}

\begin{proposition}[Digit-first sign recovery]\label{prop:digit}
Every integer zero of the new source satisfies
\[
 B=\binom{F^2+t}{t}>0,\qquad r\ge0,\qquad s\ge0,\qquad
 F\ge1,\qquad B\ge2,\qquad o\ge0.
\]
\end{proposition}
\begin{proof}
The factorial numerator is exactly $U_f=(q+1)^n$.
Let
\[
 \rho=\sum_{j=0}^{t-1}\binom njq^j,\qquad d=\binom nt.
\]
The sum of all binomial coefficients is $2^n$. Since the omitted
coefficients include a positive one, Lemma~\ref{lem:coefficient} gives
\begin{align*}
 0\le\rho
   &\le(2^n-1)q^{t-1}<q^t/2,\\
 0<d&<2^n\le q/2.
\end{align*}
In particular $\rho$ is the ordinary low remainder of $U_f$ modulo
$q^t$. Since $v_t=M/q^t>0$, the canonical normalized low remainder
$v_t\rho$ lies in $[0,M/2)$. Equation \eqref{eq:digit} gives
$R_f\equiv v_t\rho\pmod M$, whereas \eqref{eq:norm} gives
$|R_f|<M/2$. Their difference is a multiple of $M$ of absolute
value less than $M$, so $R_f=v_t\rho$.

After subtraction and division by $M$, the digit row implies
$B\equiv d\pmod q$. Lemma~\ref{lem:radix} gives $|B|<q/2$,
and $0<d<q/2$. Hence $|B-d|<q$ and the congruence is equality.
This establishes the positive digit before the product guard has
been used.

Now $rs=\G(\alpha,\beta,\gamma)\ge0$ and $r+s=B-1\ge0$ imply
$r,s\ge0$: opposite strict signs contradict the product, two
negative signs contradict the sum, and a negative sign paired with
zero also contradicts the sum. The ceiling row gives
$BF=V+s>0$, so $F\ge1$. Moreover $V=(F^2+1)^t>F^2$,
whence $BF>F^2$ and $B>F\ge1$.
Finally $q=(BH_0)^x\ge B^x>B$, so
$o=q-B-1\ge0$.
\end{proof}

All three deleted scalar guards can now be filled by fresh ternary
representations. The order of the proof is essential: it does not
assume the ceiling signs in order to obtain the binomial digit that
establishes those signs. The retained ratio for $q^{t+1}$, previously
serving the common-remainder normalization, supplies the early
information that breaks this dependence.

\section{Returning to the calibrated source}
\label{sec:factorial}

For completeness of the mathematical exposition, we give the factorial
fixed-point step explicitly. The rest of the CRT and Horn-flow
decoding is retained without alteration.

\begin{lemma}[The square-output factorial fixed point]\label{lem:factorial}
For the positive values recovered above, $F=t!$.
\end{lemma}
\begin{proof}
Write $T=t!$. From $r=B-s-1\ge0$ and $s\ge0$ we have
$0\le s\le B-1$, hence
\[
 B(F-1)<V\le BF,\qquad F=\left\lceil V/B\right\rceil.
\]
The recovered binomial identity gives
\[
 \frac VB
  =T\prod_{i=1}^{t}\frac{F^2+1}{F^2+i}\le T.
\]
As $T$ is integral, $F\le T$. Each factor in the product is
nondecreasing as a function of $F^2\ge1$. Evaluating at $F^2=1$
therefore gives
\[
 \frac VB\ge \frac{2^t}{t+1}>t^2\qquad(t\ge16).
\]
The final inequality holds at $16$ and is preserved upon increasing
$t$: $2t^2>(t+1)(t+2)$ for $t\ge4$.
Thus $F>t^2$.

For $0\le e_i<1$, the elementary product inequality
$\prod(1-e_i)\ge1-\sum e_i$ now gives
\[
 \prod_{i=1}^{t}\frac{F^2+1}{F^2+i}
 =\prod_{i=1}^{t}\left(1-\frac{i-1}{F^2+i}\right)
 >1-\frac{t(t-1)}{2F^2}
 >1-\frac{1}{2F}.
\]
Consequently
$T(1-1/(2F))<V/B\le F$, so
\[
 T<\frac{2F^2}{2F-1}
   =F+\frac{F}{2F-1}<F+1.
\]
Since $F,T$ are integers, $T\le F$. Together with $F\le T$
this proves equality.
\end{proof}

For $1\le i\le t$, put $d_i=1+F^2/i$. These are integers and
$B=\prod_i d_i$. Every prime factor of $d_i$ exceeds $t$:
if $p\le t$, then $p\mid F^2/i$, since $F=t!$, so $p\nmid d_i$.
A common prime divisor of $d_i,d_j$ would divide
$id_i-jd_j=i-j$, impossible when $i\ne j$. Thus the $d_i$ are
pairwise coprime. These are the inherited amplified CRT moduli.

After the front signs and factorial have been recovered, the proof in
\file{source/inherited_common/proofs/NORMALIZED_REMAINDERS_PROOF.md}
applies to the remaining ten low remainders. It covers the zero
representative case using
$\binom aj\le\max(1,a)^k$, rather than an invalid bound by $0^k$.
For each relevant denominator $L_i\mid M$, its genuine remainder
$r_i^0$ satisfies $0\le r_i^0<L_i/32$.
The norm and congruence then force $R_i=(M/L_i)r_i^0$ exactly.
This restores all old quotient and interval meanings.

The unchanged prime-power range tests and unit-weight Horn-flow
decoding give a finite C222 derivation of the contradiction code.
The retained calibration supplies the two links
\[
 \Ccal\text{ derives contradiction}
 \ \longleftrightarrow\
 \ZF\text{ has a finite refutation}
 \ \longleftrightarrow\
 \ZFC\text{ has a finite refutation}.
\]
Their named endpoints are
\file{KEndpoint.c222_iff_zf_refutation} and
\file{ConstructibleChoice.zf_refutation_iff_zfc_refutation}.
The corresponding Lean source files and the inherited ordinary
calibration documents are included under
\file{source/inherited_common/upstream/round17/architecture/proofs/inherited/calibration/}.
This proves soundness of the new source.

\section{Completeness of the square restriction}
\label{sec:complete}

It is not enough to observe that \eqref{eq:new} is sound. The new
restriction $q=(6\zeta+3)^2$ must still allow a witness for every
finite refutation.

Begin with a finite refutation and the predecessor's cofinal arithmetic
construction. Choose a permitted horizon $t=16+h_0^2$ large enough,
put $F=t!$, $V=(F^2+1)^t$, and $y=V(V-1)$. The beta-CRT
moduli $1+(i+1)y$ are all odd. A CRT residue class modulo their
odd least common multiple therefore contains even, arbitrarily large
integers. Choose $x$ even and sufficiently large to satisfy all the
retained endpoint bounds.

Since $i\mid F$ and $6\mid F$, each $F^2/i=F(F/i)$ is divisible
by $6$. Thus
\begin{equation}\label{eq:modsix}
 B=\prod_{i=1}^{t}(1+F^2/i)\equiv1\pmod6,\qquad
 N=B^x\equiv1\pmod6,\qquad N^2\equiv1\pmod6.
\end{equation}
The independent representative moduli have product $N^2$.
Adding a multiple of $N^2$ to any actual representative preserves
all its encoded local data and permits an arbitrary residue modulo $6$.

One must preserve the semantic anchor while making this adjustment.
Choose the ordinary $b$ representative first. If $a_0$ is the least
nonnegative residue for the dominant $a$ representative, take
\[
 a=a_0+8N^2\quad\hbox{or}\quad a=a_0+9N^2
\]
to make $a,b$ have opposite parities. The literal derived-label
constructor has $h\ge2a^2$ and $h\equiv1+a+b\pmod2$.
Hence $h$ is even and at least $128N^4$.

Reserve one ordinary representative other than $a,b$. Start it in
$(0,N^2]$, then add $jN^2$ for an appropriate
$j\in\{0,\ldots,5\}$ to arrange
\begin{equation}\label{eq:Hthree}
 H_0=1+\sum_i a_i\equiv3\pmod6.
\end{equation}
The adjusted representative remains at most $6N^2$; the dominant
$a$ remains below $10N^2$. The other representatives can remain
below $2N^2$, or $6N^2$ if harmless earlier parity choices are
retained. The fixed front indices are below $N^2$, as in the parent.
Thus the even $h\ge128N^4$ still dominates every other requested
index, and the squared-cap completeness estimates remain valid.
The actual radix and the private frame are built after these lifts;
there is no attempt to keep an obsolete upper bound on $H_0$.

By \eqref{eq:modsix} and \eqref{eq:Hthree}, $BH_0$ is an odd
multiple of $3$. As $x$ is positive and even,
$Q=(BH_0)^{x/2}$ is also an odd multiple of $3$. Therefore
\[
 \zeta=(Q-3)/6\in\ZZ,\qquad q=(BH_0)^x=(6\zeta+3)^2.
\]
No fixed $3$-adic valuation is being imposed; every odd multiple of
$3$ is congruent to $3$ modulo $6$.

The canonical ceiling remainders $r,s$ are nonnegative, so their
product has a $\G$ representation. The eleven genuine normalized
remainders satisfy $0\le R_i<M/32$. Hence
\[
 4\sum_{i=1}^{11}R_i^2<\frac{11}{256}M^2<M^2.
\]
The integer $M^2-1-4\sum R_i^2$ is nonnegative and has its own
independent $\G$ representation, supplying \eqref{eq:norm}.
Finally choose the frame and the canonical private Pell witnesses,
including the inherited deterministic comparison center. All new
source rows are satisfied.

This is an existential construction with fresh private witnesses.
It does not claim a bijection of all old and new solution tuples.
Together with Section~\ref{sec:factorial}, it proves both directions
of the source equivalence.

\section{From the source theorem to the literal quartic}
\label{sec:compile}

The source compiler and certificates are included so that the displayed
equation is a finite object with a reviewable connection to the proof.
No search for witnesses or assumption of consistency is used to
construct its coefficients.

\subsection{Exact algebraic transformations}

The file \file{source/equation247_weak/delta_certificate.json}
records every changed source row and every deleted private guard.
The source audit checks the actual polynomials, not just row names.
Eighteen inherited Horn factorizations and 204 integral monic
restorations lead to the raw 247-variable system; seventy integral
translations give its centered coordinates.

A monic restoration deletes a variable only when an explicit polynomial
with integer coefficients reconstructs it from the retained variables.
The reverse substitution is checked exactly. Thus there are no
unreported divisibility conditions or rational coordinates.
Unused private coordinates are omitted only after their nonoccurrence
has been verified.

The final coordinate census explains the net saving precisely.
The nine old scalar-guard coordinates and the old free center coordinate
disappear. Four fresh guard coordinates appear, and the coordinate
\file{fact0.remainder}, which was eliminable in the predecessor, is
retained to keep the new product row quadratic. Thus the measured
change is $252-10+5=247$, rather than a naive count of deleted guards.

To improve the final coefficient profile, we transfer 76 known
integral coordinate shears from the prior shared-residue packet
\cite{shared}, matching the original coordinate names rather than
assuming the same coordinate numbering. A step $(i,j,k)$ means
\[
 \text{old }z_j=\text{new }z_j+k\,\text{new }z_i.
\]
It has determinant one and inverse coefficient $-k$.
The entire integer matrix, its two-sided inverse, and the full old-to-new
coordinate map are checked. These are coordinate changes on
$\ZZ^{247}$, so integer solvability is preserved.

Four additional residual operations improve the height:
\[
 R_{18}\leftarrow R_{18}-22R_{11},\quad
 R_{19}\leftarrow R_{19}-R_{11},\quad
 R_{52}\leftarrow R_{52}-R_{11},\quad
 R_{27}\leftarrow R_{27}-R_{23}.
\]
The indices here are one-based and the operations occur in the displayed
order. They are recorded with zero-based indices in
\file{certificates/row_operations.json}.
Each operation is an invertible integer row operation. It preserves
the simultaneous zero set of the residuals. It generally changes
the scalar sum-of-squares polynomial, so we do not assert an identity
between the old and new sums of squares.

All final residuals remain quadratic, with maximum absolute residual
coefficient $66$ and 1,077 total residual monomial occurrences.
Their squared sum has the exact metrics in Section~\ref{sec:equation};
all 247 variables occur in the collected polynomial.

\subsection{What has been checked}

\file{source/verify_composed.py} independently reconstructs the source
delta, invokes the retained independently implemented Horn and monic
replay, checks translations, and collects the raw quartic.
\file{code/build_final.py} reconstructs the final finite outputs from
the frozen coordinate and residual certificates.

\file{code/check_final.py} is a second implementation using sparse
exponent-vector monomials and explicit binomial expansion. It imports
no producer or eliminator. It checks each map forward and backward,
both matrix inverse identities, every final residual, every collected
coefficient, the complete coordinate map, degree, height, support,
and coordinate liveness. It also parses the actual LaTeX residual
file back into polynomials and compares them with the machine-readable
equation.

Six deliberate source/projection corruptions and eight deliberate
final-algebra corruptions are rejected. The latter include changes
to a shear, its historical coordinate mapping, a residual operation,
a final residual, a collected coefficient, the coordinate map,
a printed formula, and the domain. Normal Python and optimized
Python both pass; disabling assertions does not disable the checks.
The receipts and substantive reviews are retained in \file{review/}.

These finite checks prove exact algebraic identities for the emitted
finite data. They do not replace the unbounded number-theoretic
arguments or recheck the entire inherited Lean development.
The proof of Theorem~\ref{thm:main} is the ordinary argument in
Sections~\ref{sec:interface}--\ref{sec:complete}, composed with the
retained calibration and the verified finite transformations.
Every new argument uses elementary arithmetic, finite coding and
standard set-theoretic constructions formalizable in ZFC. No fresh
kernel-checked proof of this complete final theorem, and no encoded
object-ZFC proof term, is claimed.

\section{Reproducibility, attribution, and further research}
\label{sec:research}

The canonical starting point is \file{README.md}. With Python's
standard library, run
\begin{verbatim}
python3 -B verify.py
python3 -O -B verify.py
\end{verbatim}
The default checks are read-only. To reconstruct the final files from
the frozen raw source and maps, run
\begin{verbatim}
python3 -B code/build_final.py --write
\end{verbatim}
To evaluate an explicitly supplied 247-integer tuple without pretending
to decide consistency, run
\begin{verbatim}
python3 -B code/evaluate.py tuple.json --expanded
\end{verbatim}
The input is a JSON list in the order $z_1,\ldots,z_{247}$.
The build script \file{build_report.sh} compiles this report and the
optional experimental note. The package includes the ordinary source
proofs, exact transformation certificates, independent checkers,
machine-readable equations, all report inputs, a provenance ledger,
and a SHA-256 manifest.

\subsection{The actual new contribution}

The deterministic Pell center, odd-square idea, common normalized
remainder norm, and historical coordinate shears are credited to
their incoming source packets \cite{center,budgets,common,shared}.
The new mathematical contribution here is the specific
\emph{digit-first sign recovery} argument:
\[
 q^{t+1}\mid q^x
 \ \Longrightarrow\ x\ge t+1
 \ \Longrightarrow\ 2^{F^2+t}\le q/2
 \ \Longrightarrow\ B=\binom{F^2+t}{t}>0.
\]
The self-radix bound on $|B|$ is also essential in the last implication.
This order permits the weak product guard and removes the three
individual front sign certificates. The nonunit square
$(6\zeta+3)^2$ and its cofinal CRT construction make the early radix
argument unconditional.

This is a concrete improvement relative to the audited snapshot.
It is not an exhaustive literature-priority claim or a declaration
of a major breakthrough. General universal-equation results
\cite{jones,bayer,sun} establish useful context for variable-count
tradeoffs but do not supply this direct source theorem.

\subsection{A different calculus was implemented and measured}

The accompanying directory
\file{experiments/three_parameter_calculus/} investigates replacing
the four-parameter implicational schema by K, S, and Peirce, together
with a modular saturation equation whose inverse coordinate is
unbounded. A literal 521-step Hilbert derivation certifies the
propositional replacement. Its note proves the required local
small-value decoding, prime-divisor argument, CRT inverse choice,
and nonnegative quotient repair.

The resulting experimental quartic has 231 \emph{nonnegative}
integer variables, 183 residuals, 6,369 monomials and height
83,899,166. It loses to its own 229-variable nonnegative parent:
extra source branches and multiplication gates cost more than the
removed bounded column saves. These domains and architectures
must not be confused with the 247 unrestricted integers of
Theorem~\ref{thm:main}. The experiment is included because it
provides an exact, falsifiable result for a distinct proposed route.
Its standalone checker verifies the finite basis certificate and
polynomial expansion; rebuilding through its different full parent
compiler additionally requires the pinned parent archive identified
in that directory.

\subsection{Questions suggested by the reduction}

The most promising next work is to examine existing divisibility
identities for information available before the sign bootstrap.
The $q^{t+1}$ ratio is indispensable to the present proof; deleting
it on the assumption that it is merely a remainder-normalization
overhead would destroy the early exponent bound.
Other ratios may have a different status, but a reduction must
preserve both quadratic degree and an integral restoration.

One can also ask whether the single product nonnegativity condition
can be represented with fewer coordinates under additional cofinal
arithmetic restrictions. Such a claim needs a complete cofinal
construction; a sign argument alone is insufficient.
The unsuccessful calculus experiment suggests co-designing the
local proof calculus with its arithmetic circuit, since fewer
logical parameters need not mean fewer final polynomial variables.

Finally, the integer coordinate and residual searches can be improved
independently of the semantics, provided every chosen map has a
checked integral inverse. Formalizing the early-exponent and
digit-sign lemmas in the existing proof environment would make a
valuable next integration step. None of these proposed directions
is used as an unproved premise in the delivered equation.

\begin{thebibliography}{9}
\bibitem{repo}
V. Reshetnikov et al.,
\emph{ProveIt: Computability and Diophantine consistency research},
snapshot \file{3ecd6a6cf53d47519228328567f45fb7ee9781d9}.
\url{https://github.com/VladimirReshetnikov/ProveIt/tree/3ecd6a6cf53d47519228328567f45fb7ee9781d9/Computability/DiophantineConsistency/con-zfc-diophantine-forms}.
The ledger also pins the incoming archives.

\bibitem{common}
\emph{ConZFC Common Remainders / Radix 47}, 25 October 2026,
\file{docs/incoming/ConZFC_Common_Remainders_Radix47_2026-10-25.zip},
at snapshot \cite{repo}. The portable direct source closure and its
ordinary proofs are retained under \file{source/inherited_common/}.

\bibitem{center}
\emph{ConZFC Symmetric Remainders}, 25 October 2026,
\file{docs/incoming/ConZFC_Symmetric_Remainders_2026-10-25.zip},
at snapshot \cite{repo}. Deterministic comparison-center proof
retained as \file{proofs/PELL_CENTER_PRIOR.md}.

\bibitem{budgets}
\emph{ConZFC Symmetric Residue Budgets}, 25 October 2026,
\file{docs/incoming/ConZFC_Symmetric_Residue_Budgets_2026-10-25.zip},
at snapshot \cite{repo}. Prior odd-square proof retained as
\file{proofs/ODD_SQUARE_PRIOR.md}.

\bibitem{shared}
\emph{ConZFC Shared Residue Bounds}, 25 October 2026,
\file{docs/incoming/ConZFC_Shared_Residue_Bounds_2026-10-25.zip},
at snapshot \cite{repo}. Exact historical shear inputs and transfer
certificate are retained in \file{certificates/}.

\bibitem{jones}
J. P. Jones, Universal Diophantine equation,
\emph{Journal of Symbolic Logic} 47 (1982), 549--571.
\url{https://doi.org/10.2307/2273588}.

\bibitem{bayer}
J. Bayer, M. David, M. Hassler, Y. Matiyasevich and D. Schleicher,
\emph{Diophantine Equations over $\ZZ$: Universal Bounds and
Parallel Formalization}, arXiv:2506.20909, version 3, 22 October 2026.
\url{https://arxiv.org/abs/2506.20909}.

\bibitem{sun}
Z.-W. Sun, \emph{Ten unknowns for Hilbert's tenth problem over the integers},
arXiv:2609.01594, version 2, 2 September 2026.
\url{https://arxiv.org/abs/2609.01594}.
\end{thebibliography}

\clearpage
\appendix
\section{All 182 quadratic residuals}\label{app:rows}

Every symbol below is one of the 247 unrestricted integer coordinates.
The single equation is the sum of the squares of these complete
residuals, set equal to zero. Coefficients omitted before a monomial
are $1$ or $-1$ according to the displayed sign. There are no
additional definitions or constraints in this appendix.

\begingroup
\small
\setlength{\abovedisplayskip}{5pt}
\setlength{\belowdisplayskip}{5pt}
\setlength{\abovedisplayshortskip}{3pt}
\setlength{\belowdisplayshortskip}{3pt}
% Complete final quadratic residuals. Generated by code/build_final.py.
\begin{equation*}\begin{aligned}
R_{1}&=-z_{4}^{2} +z_{5}
\end{aligned}\end{equation*}
\begin{equation*}\begin{aligned}
R_{2}&=-1 -z_{4}\,z_{6} +z_{6} +z_{7} \\
&\quad -z_{8}
\end{aligned}\end{equation*}
\begin{equation*}\begin{aligned}
R_{3}&=z_{1} -z_{1}\,z_{3} -z_{3} +z_{11}
\end{aligned}\end{equation*}
\begin{equation*}\begin{aligned}
R_{4}&=-1 +z_{2} -z_{3}\,z_{9} -z_{4}
\end{aligned}\end{equation*}
\begin{equation*}\begin{aligned}
R_{5}&=-1 +z_{2} -z_{10}\,z_{11}
\end{aligned}\end{equation*}
\begin{equation*}\begin{aligned}
R_{6}&=1 +z_{5} -z_{6}\,z_{22} +z_{20}
\end{aligned}\end{equation*}
\begin{equation*}\begin{aligned}
R_{7}&=-1 +z_{12} -z_{12}^{2} -2\,z_{12}\,z_{13} \\
&\quad -z_{13} -z_{13}^{2} +z_{23}
\end{aligned}\end{equation*}
\begin{equation*}\begin{aligned}
R_{8}&=-5 -20\,z_{12} -16\,z_{12}^{2} -16\,z_{12}\,z_{13} \\
&\quad -16\,z_{12}\,z_{23} -10\,z_{13} -4\,z_{13}^{2} -8\,z_{13}\,z_{23} \\
&\quad -10\,z_{23} -4\,z_{23}^{2} +z_{24}
\end{aligned}\end{equation*}
\begin{equation*}\begin{aligned}
R_{9}&=-2 -z_{12} -z_{12}^{2} -4\,z_{12}\,z_{24} \\
&\quad -4\,z_{24} -4\,z_{24}^{2} +z_{25}
\end{aligned}\end{equation*}
\begin{equation*}\begin{aligned}
R_{10}&=1 +2\,z_{25} -4\,z_{25}^{2} +z_{26}
\end{aligned}\end{equation*}
\begin{equation*}\begin{aligned}
R_{11}&=-2 -z_{13}^{2} -2\,z_{13}\,z_{14} -z_{14} \\
&\quad -z_{14}^{2} +z_{27}
\end{aligned}\end{equation*}
\begin{equation*}\begin{aligned}
R_{12}&=-9 +7\,z_{12} -z_{12}^{2} -4\,z_{12}\,z_{13} \\
&\quad +4\,z_{12}\,z_{14} -4\,z_{12}\,z_{27} +12\,z_{13} -4\,z_{13}^{2} \\
&\quad +8\,z_{13}\,z_{14} -8\,z_{13}\,z_{27} -15\,z_{14} -4\,z_{14}^{2} \\
&\quad +8\,z_{14}\,z_{27} +12\,z_{27} -4\,z_{27}^{2} +z_{28} \\
&\quad -z_{34}
\end{aligned}\end{equation*}
\begin{equation*}\begin{aligned}
R_{13}&=-1 +z_{12} -z_{12}^{2} -2\,z_{12}\,z_{14} \\
&\quad -z_{13} -z_{14}^{2} +z_{29}
\end{aligned}\end{equation*}
\begin{equation*}\begin{aligned}
R_{14}&=-4\,z_{12} -64\,z_{12}^{2} -64\,z_{12}\,z_{14} -32\,z_{12}\,z_{23} \\
&\quad -32\,z_{12}\,z_{29} -2\,z_{13} -16\,z_{14}^{2} -16\,z_{14}\,z_{23} \\
&\quad -16\,z_{14}\,z_{29} -2\,z_{23} -4\,z_{23}^{2} -8\,z_{23}\,z_{29} \\
&\quad -z_{28} -4\,z_{29}^{2} +z_{30} +z_{34}
\end{aligned}\end{equation*}
\begin{equation*}\begin{aligned}
R_{15}&=-36\,z_{14}^{2} +24\,z_{14}\,z_{30} -4\,z_{30}^{2} +z_{31}
\end{aligned}\end{equation*}
\begin{equation*}\begin{aligned}
R_{16}&=-6 -8\,z_{12} -4\,z_{12}^{2} -8\,z_{12}\,z_{13} \\
&\quad -8\,z_{12}\,z_{14} +8\,z_{12}\,z_{27} +8\,z_{12}\,z_{29} -6\,z_{13} \\
&\quad -4\,z_{13}^{2} -8\,z_{13}\,z_{14} +8\,z_{13}\,z_{27} +8\,z_{13}\,z_{29} \\
&\quad -5\,z_{14} -4\,z_{14}^{2} +8\,z_{14}\,z_{27} +8\,z_{14}\,z_{29} \\
&\quad -z_{23} +8\,z_{27} -4\,z_{27}^{2} -8\,z_{27}\,z_{29} \\
&\quad +6\,z_{29} -4\,z_{29}^{2} +z_{32}
\end{aligned}\end{equation*}
\begin{equation*}\begin{aligned}
R_{17}&=1 -4\,z_{12} +2\,z_{13} -4\,z_{13}^{2} \\
&\quad -8\,z_{13}\,z_{14} -8\,z_{13}\,z_{32} +4\,z_{14} -4\,z_{14}^{2} \\
&\quad -8\,z_{14}\,z_{32} -2\,z_{23} +4\,z_{32} -4\,z_{32}^{2} \\
&\quad +z_{33}
\end{aligned}\end{equation*}
\begin{equation*}\begin{aligned}
R_{18}&=34 -5\,z_{12} -z_{12}^{2} -12\,z_{12}\,z_{13} \\
&\quad -12\,z_{12}\,z_{14} -4\,z_{12}\,z_{27} -36\,z_{13} -14\,z_{13}^{2} \\
&\quad -28\,z_{13}\,z_{14} -24\,z_{13}\,z_{27} -11\,z_{14} -14\,z_{14}^{2} \\
&\quad -24\,z_{14}\,z_{27} -34\,z_{27} -4\,z_{27}^{2} +z_{34}
\end{aligned}\end{equation*}
\begin{equation*}\begin{aligned}
R_{19}&=-16\,z_{12} -16\,z_{12}^{2} +8\,z_{12}\,z_{13} -32\,z_{12}\,z_{14} \\
&\quad -16\,z_{12}\,z_{29} +5\,z_{13} +10\,z_{13}\,z_{14} +4\,z_{13}\,z_{29} \\
&\quad -14\,z_{14} -15\,z_{14}^{2} -16\,z_{14}\,z_{29} -z_{27} \\
&\quad -8\,z_{29} -4\,z_{29}^{2} -z_{34} +z_{35}
\end{aligned}\end{equation*}
\begin{equation*}\begin{aligned}
R_{20}&=2 -64\,z_{14}^{2} -32\,z_{14}\,z_{35} -4\,z_{35}^{2} \\
&\quad +z_{36}
\end{aligned}\end{equation*}
\begin{equation*}\begin{aligned}
R_{21}&=-12 -3\,z_{12} -47\,z_{13} -36\,z_{13}^{2} \\
&\quad +24\,z_{13}\,z_{23} +24\,z_{13}\,z_{27} -z_{14} +16\,z_{23} \\
&\quad -4\,z_{23}^{2} -8\,z_{23}\,z_{27} +14\,z_{27} -4\,z_{27}^{2} \\
&\quad -z_{29} +z_{37}
\end{aligned}\end{equation*}
\begin{equation*}\begin{aligned}
R_{22}&=-3 -6\,z_{12} +2\,z_{13} -2\,z_{14} \\
&\quad -2\,z_{29} -4\,z_{37} -4\,z_{37}^{2} +z_{38}
\end{aligned}\end{equation*}
\begin{equation*}\begin{aligned}
R_{23}&=-1 -13\,z_{12} -16\,z_{12}^{2} -16\,z_{12}\,z_{13} \\
&\quad -16\,z_{12}\,z_{23} -6\,z_{13} -4\,z_{13}^{2} -8\,z_{13}\,z_{23} \\
&\quad -6\,z_{23} -4\,z_{23}^{2} +z_{39}
\end{aligned}\end{equation*}
\begin{equation*}\begin{aligned}
R_{24}&=-2 +3\,z_{12} -z_{12}^{2} +4\,z_{12}\,z_{39} \\
&\quad -4\,z_{39} -4\,z_{39}^{2} +z_{40}
\end{aligned}\end{equation*}
\begin{equation*}\begin{aligned}
R_{25}&=1 +2\,z_{40} -4\,z_{40}^{2} +z_{41}
\end{aligned}\end{equation*}
\begin{equation*}\begin{aligned}
R_{26}&=-z_{13}^{2} -4\,z_{13}\,z_{41} -4\,z_{41}^{2} +z_{42}
\end{aligned}\end{equation*}
\begin{equation*}\begin{aligned}
R_{27}&=-5 -3\,z_{12} +48\,z_{12}\,z_{13} -16\,z_{12}\,z_{23} \\
&\quad +12\,z_{13} -12\,z_{13}^{2} +40\,z_{13}\,z_{23} -2\,z_{14} \\
&\quad -8\,z_{23} -12\,z_{23}^{2} -z_{27} -z_{39} \\
&\quad +z_{43}
\end{aligned}\end{equation*}
\begin{equation*}\begin{aligned}
R_{28}&=-4 -3\,z_{12} -7\,z_{13} -4\,z_{14} \\
&\quad -2\,z_{27} -z_{29} +4\,z_{43} -4\,z_{43}^{2} \\
&\quad +z_{44}
\end{aligned}\end{equation*}
\begin{equation*}\begin{aligned}
R_{29}&=-6\,z_{12} +2\,z_{13} +2\,z_{14} -4\,z_{14}^{2} \\
&\quad -8\,z_{14}\,z_{44} -2\,z_{29} +4\,z_{44} -4\,z_{44}^{2} \\
&\quad +z_{45}
\end{aligned}\end{equation*}
\begin{equation*}\begin{aligned}
R_{30}&=-z_{14} -z_{14}^{2} -2\,z_{14}\,z_{15} -z_{15}^{2} \\
&\quad +z_{46}
\end{aligned}\end{equation*}
\begin{equation*}\begin{aligned}
R_{31}&=2\,z_{12} -z_{12}^{2} -4\,z_{12}\,z_{46} +2\,z_{46} \\
&\quad -4\,z_{46}^{2} +z_{47} -z_{48}
\end{aligned}\end{equation*}
\begin{equation*}\begin{aligned}
R_{32}&=-35 +2\,z_{13} -66\,z_{14} -36\,z_{14}^{2} \\
&\quad +24\,z_{14}\,z_{27} +24\,z_{14}\,z_{29} +24\,z_{27} -4\,z_{27}^{2} \\
&\quad -8\,z_{27}\,z_{29} +22\,z_{29} -4\,z_{29}^{2} +z_{48}
\end{aligned}\end{equation*}
\begin{equation*}\begin{aligned}
R_{33}&=-16\,z_{14}^{2} -16\,z_{14}\,z_{47} -2\,z_{47} -4\,z_{47}^{2} \\
&\quad +2\,z_{48} +z_{49}
\end{aligned}\end{equation*}
\begin{equation*}\begin{aligned}
R_{34}&=-16\,z_{12}^{2} +24\,z_{12}\,z_{14} -16\,z_{12}\,z_{23} +16\,z_{12}\,z_{32} \\
&\quad +z_{14} -9\,z_{14}^{2} +12\,z_{14}\,z_{23} -12\,z_{14}\,z_{32} \\
&\quad -4\,z_{23}^{2} +8\,z_{23}\,z_{32} -4\,z_{32}^{2} +z_{50}
\end{aligned}\end{equation*}
\begin{equation*}\begin{aligned}
R_{35}&=-1 +6\,z_{13} -4\,z_{13}^{2} +8\,z_{13}\,z_{14} \\
&\quad -8\,z_{13}\,z_{27} -6\,z_{14} -4\,z_{14}^{2} +8\,z_{14}\,z_{27} \\
&\quad +6\,z_{27} -4\,z_{27}^{2} +z_{51}
\end{aligned}\end{equation*}
\begin{equation*}\begin{aligned}
R_{36}&=-1 -z_{14} -z_{14}^{2} -4\,z_{14}\,z_{51} \\
&\quad -4\,z_{51}^{2} +z_{52}
\end{aligned}\end{equation*}
\begin{equation*}\begin{aligned}
R_{37}&=1 +z_{13} -z_{13}^{2} -4\,z_{13}\,z_{52} \\
&\quad -4\,z_{52}^{2} +z_{53}
\end{aligned}\end{equation*}
\begin{equation*}\begin{aligned}
R_{38}&=-1 -6\,z_{53} -4\,z_{53}^{2} +z_{54}
\end{aligned}\end{equation*}
\begin{equation*}\begin{aligned}
R_{39}&=-1 -z_{14} -z_{14}^{2} -4\,z_{14}\,z_{54} \\
&\quad -4\,z_{54}^{2} +z_{55}
\end{aligned}\end{equation*}
\begin{equation*}\begin{aligned}
R_{40}&=2 +z_{13} -z_{13}^{2} -4\,z_{13}\,z_{55} \\
&\quad -4\,z_{55}^{2} +z_{56}
\end{aligned}\end{equation*}
\begin{equation*}\begin{aligned}
R_{41}&=-2\,z_{50} -4\,z_{50}^{2} -8\,z_{50}\,z_{56} -4\,z_{56}^{2} \\
&\quad +z_{57}
\end{aligned}\end{equation*}
\begin{equation*}\begin{aligned}
R_{42}&=-9 +7\,z_{12} -z_{12}^{2} -4\,z_{12}\,z_{57} \\
&\quad +12\,z_{57} -4\,z_{57}^{2} +z_{58}
\end{aligned}\end{equation*}
\begin{equation*}\begin{aligned}
R_{43}&=-1 -6\,z_{58} -4\,z_{58}^{2} +z_{59}
\end{aligned}\end{equation*}
\begin{equation*}\begin{aligned}
R_{44}&=-z_{14} -z_{14}^{2} -4\,z_{14}\,z_{59} -4\,z_{59}^{2} \\
&\quad +z_{60}
\end{aligned}\end{equation*}
\begin{equation*}\begin{aligned}
R_{45}&=-z_{12} -z_{12}^{2} -4\,z_{12}\,z_{60} -4\,z_{60} \\
&\quad -4\,z_{60}^{2} +z_{61}
\end{aligned}\end{equation*}
\begin{equation*}\begin{aligned}
R_{46}&=-10 -z_{13} +z_{14} -z_{27} \\
&\quad -14\,z_{61} -4\,z_{61}^{2} +z_{62}
\end{aligned}\end{equation*}
\begin{equation*}\begin{aligned}
R_{47}&=1 -2\,z_{13} +2\,z_{14} -2\,z_{27} \\
&\quad +4\,z_{62} -4\,z_{62}^{2} +z_{63}
\end{aligned}\end{equation*}
\begin{equation*}\begin{aligned}
R_{48}&=z_{13} -z_{13}^{2} -4\,z_{13}\,z_{63} -4\,z_{63}^{2} \\
&\quad +z_{64} -z_{65}
\end{aligned}\end{equation*}
\begin{equation*}\begin{aligned}
R_{49}&=-10 -14\,z_{56} -4\,z_{56}^{2} +z_{65}
\end{aligned}\end{equation*}
\begin{equation*}\begin{aligned}
R_{50}&=1 -2\,z_{64} -4\,z_{64}^{2} +2\,z_{65} \\
&\quad +z_{66}
\end{aligned}\end{equation*}
\begin{equation*}\begin{aligned}
R_{51}&=z_{14} -z_{14}^{2} -4\,z_{14}\,z_{66} -4\,z_{66}^{2} \\
&\quad +z_{67}
\end{aligned}\end{equation*}
\begin{equation*}\begin{aligned}
R_{52}&=2 +15\,z_{13} -35\,z_{13}^{2} +14\,z_{13}\,z_{14} \\
&\quad +24\,z_{13}\,z_{23} -4\,z_{14}\,z_{23} -5\,z_{23} -4\,z_{23}^{2} \\
&\quad -z_{27} +z_{68}
\end{aligned}\end{equation*}
\begin{equation*}\begin{aligned}
R_{53}&=1 +6\,z_{13} -2\,z_{23} -z_{67} \\
&\quad +4\,z_{68} -4\,z_{68}^{2} +z_{69}
\end{aligned}\end{equation*}
\begin{equation*}\begin{aligned}
R_{54}&=-2 +2\,z_{67} -2\,z_{69} -4\,z_{69}^{2} \\
&\quad +z_{70}
\end{aligned}\end{equation*}
\begin{equation*}\begin{aligned}
R_{55}&=-2 -4\,z_{70} -4\,z_{70}^{2} +z_{71}
\end{aligned}\end{equation*}
\begin{equation*}\begin{aligned}
R_{56}&=-z_{4}\,z_{17} +z_{72}
\end{aligned}\end{equation*}
\begin{equation*}\begin{aligned}
R_{57}&=z_{16} -z_{17} -2\,z_{26}\,z_{72} -z_{73}
\end{aligned}\end{equation*}
\begin{equation*}\begin{aligned}
R_{58}&=12\,z_{14}\,z_{72} +z_{16} -z_{17} -4\,z_{28}\,z_{72} \\
&\quad -2\,z_{31}\,z_{72} +4\,z_{34}\,z_{72} -z_{74}
\end{aligned}\end{equation*}
\begin{equation*}\begin{aligned}
R_{59}&=z_{16} -z_{17} -2\,z_{33}\,z_{72} -6\,z_{72} \\
&\quad -z_{75}
\end{aligned}\end{equation*}
\begin{equation*}\begin{aligned}
R_{60}&=-12\,z_{14}\,z_{72} +z_{16} -z_{17} -4\,z_{34}\,z_{72} \\
&\quad -2\,z_{36}\,z_{72} -z_{76}
\end{aligned}\end{equation*}
\begin{equation*}\begin{aligned}
R_{61}&=z_{16} -z_{17} -2\,z_{38}\,z_{72} +6\,z_{72} \\
&\quad -z_{77}
\end{aligned}\end{equation*}
\begin{equation*}\begin{aligned}
R_{62}&=z_{16} -z_{17} -4\,z_{41}\,z_{72} -2\,z_{42}\,z_{72} \\
&\quad -z_{78}
\end{aligned}\end{equation*}
\begin{equation*}\begin{aligned}
R_{63}&=z_{16} -z_{17} -2\,z_{45}\,z_{72} -z_{79}
\end{aligned}\end{equation*}
\begin{equation*}\begin{aligned}
R_{64}&=z_{16} -z_{17} -2\,z_{49}\,z_{72} -4\,z_{72} \\
&\quad -z_{80}
\end{aligned}\end{equation*}
\begin{equation*}\begin{aligned}
R_{65}&=-z_{12}\,z_{13} +z_{12}\,z_{14} +z_{13} +z_{13}\,z_{14} \\
&\quad -z_{14} -z_{14}^{2} +z_{81}
\end{aligned}\end{equation*}
\begin{equation*}\begin{aligned}
R_{66}&=1 +2\,z_{15} +z_{15}^{2} -z_{81}^{2} \\
&\quad +z_{82}
\end{aligned}\end{equation*}
\begin{equation*}\begin{aligned}
R_{67}&=-z_{82}^{2} +z_{83}
\end{aligned}\end{equation*}
\begin{equation*}\begin{aligned}
R_{68}&=-4\,z_{70}\,z_{72} -2\,z_{71}\,z_{72} +2\,z_{72} +z_{84}
\end{aligned}\end{equation*}
\begin{equation*}\begin{aligned}
R_{69}&=-z_{16}^{2} +2\,z_{16}\,z_{17} +2\,z_{16}\,z_{84} -z_{17}^{2} \\
&\quad -2\,z_{17}\,z_{84} -z_{84}^{2} +z_{85}
\end{aligned}\end{equation*}
\begin{equation*}\begin{aligned}
R_{70}&=-1 -z_{4}\,z_{13} +z_{86}
\end{aligned}\end{equation*}
\begin{equation*}\begin{aligned}
R_{71}&=-z_{17}\,z_{86} +z_{88}\,z_{89} +z_{90}
\end{aligned}\end{equation*}
\begin{equation*}\begin{aligned}
R_{72}&=-1 +z_{4} -z_{4}\,z_{12} +z_{87}
\end{aligned}\end{equation*}
\begin{equation*}\begin{aligned}
R_{73}&=-1 +2\,z_{4} +6\,z_{4}\,z_{13} -2\,z_{4}\,z_{23} \\
&\quad +z_{88}
\end{aligned}\end{equation*}
\begin{equation*}\begin{aligned}
R_{74}&=-z_{16}\,z_{87} +z_{89}
\end{aligned}\end{equation*}
\begin{equation*}\begin{aligned}
R_{75}&=-1 -4\,z_{4}\,z_{12} -2\,z_{4}\,z_{13} -2\,z_{4}\,z_{23} \\
&\quad +z_{91}
\end{aligned}\end{equation*}
\begin{equation*}\begin{aligned}
R_{76}&=z_{16}\,z_{86} -z_{17}\,z_{91} +z_{92}
\end{aligned}\end{equation*}
\begin{equation*}\begin{aligned}
R_{77}&=-z_{16}\,z_{73} +z_{17}\,z_{73} +z_{93}
\end{aligned}\end{equation*}
\begin{equation*}\begin{aligned}
R_{78}&=-z_{74}\,z_{93} +z_{94}
\end{aligned}\end{equation*}
\begin{equation*}\begin{aligned}
R_{79}&=-z_{75}\,z_{94} +z_{95}
\end{aligned}\end{equation*}
\begin{equation*}\begin{aligned}
R_{80}&=-z_{76}\,z_{95} +z_{96}
\end{aligned}\end{equation*}
\begin{equation*}\begin{aligned}
R_{81}&=-z_{77}\,z_{96} +z_{97}
\end{aligned}\end{equation*}
\begin{equation*}\begin{aligned}
R_{82}&=-z_{78}\,z_{97} +z_{98}
\end{aligned}\end{equation*}
\begin{equation*}\begin{aligned}
R_{83}&=-z_{79}\,z_{98} +z_{99}
\end{aligned}\end{equation*}
\begin{equation*}\begin{aligned}
R_{84}&=-z_{80}\,z_{99} +z_{100}
\end{aligned}\end{equation*}
\begin{equation*}\begin{aligned}
R_{85}&=-z_{90}\,z_{100} +z_{101}
\end{aligned}\end{equation*}
\begin{equation*}\begin{aligned}
R_{86}&=-z_{92}\,z_{101} +z_{102}
\end{aligned}\end{equation*}
\begin{equation*}\begin{aligned}
R_{87}&=-z_{102}^{2} +z_{103}
\end{aligned}\end{equation*}
\begin{equation*}\begin{aligned}
R_{88}&=-z_{3} -z_{3}\,z_{20} +z_{20} +z_{104}
\end{aligned}\end{equation*}
\begin{equation*}\begin{aligned}
R_{89}&=-z_{18}\,z_{104} +z_{105}
\end{aligned}\end{equation*}
\begin{equation*}\begin{aligned}
R_{90}&=-2\,z_{2}^{2} +2\,z_{2}\,z_{16} +2\,z_{2}\,z_{17} +2\,z_{2}\,z_{105} \\
&\quad +2\,z_{2}\,z_{106} -z_{16}^{2} -2\,z_{16}\,z_{105} -z_{17}^{2} \\
&\quad -2\,z_{17}\,z_{106} +z_{21}\,z_{108} -z_{83}\,z_{103} -z_{85}\,z_{103} \\
&\quad -z_{105}^{2} -z_{106}^{2}
\end{aligned}\end{equation*}
\begin{equation*}\begin{aligned}
R_{91}&=-z_{3}\,z_{19} +z_{19} -z_{19}\,z_{104} +z_{106}
\end{aligned}\end{equation*}
\begin{equation*}\begin{aligned}
R_{92}&=-5\,z_{12} +3\,z_{13} -4\,z_{23} -2\,z_{26} \\
&\quad -4\,z_{28} -2\,z_{31} -2\,z_{33} -2\,z_{36} \\
&\quad -2\,z_{38} -4\,z_{41} -2\,z_{42} -2\,z_{45} \\
&\quad -2\,z_{49} -4\,z_{70} -2\,z_{71} +z_{107}
\end{aligned}\end{equation*}
\begin{equation*}\begin{aligned}
R_{93}&=-z_{109}\,z_{112} +z_{112} +z_{114}
\end{aligned}\end{equation*}
\begin{equation*}\begin{aligned}
R_{94}&=-z_{6}\,z_{136} +z_{110}\,z_{235} -z_{113} -z_{114}\,z_{136}
\end{aligned}\end{equation*}
\begin{equation*}\begin{aligned}
R_{95}&=z_{115}\,z_{236} -z_{117}\,z_{139} -z_{118}
\end{aligned}\end{equation*}
\begin{equation*}\begin{aligned}
R_{96}&=z_{119}\,z_{236} -z_{120}\,z_{139} -z_{121}
\end{aligned}\end{equation*}
\begin{equation*}\begin{aligned}
R_{97}&=z_{122}\,z_{236} -z_{123}\,z_{139} -z_{124}
\end{aligned}\end{equation*}
\begin{equation*}\begin{aligned}
R_{98}&=z_{125}\,z_{236} -z_{126}\,z_{139} -z_{127}
\end{aligned}\end{equation*}
\begin{equation*}\begin{aligned}
R_{99}&=z_{128}\,z_{237} -z_{130}\,z_{139} -z_{131}
\end{aligned}\end{equation*}
\begin{equation*}\begin{aligned}
R_{100}&=z_{132}\,z_{237} -z_{133}\,z_{139} -z_{134}
\end{aligned}\end{equation*}
\begin{equation*}\begin{aligned}
R_{101}&=-z_{21}\,z_{136} +z_{139}
\end{aligned}\end{equation*}
\begin{equation*}\begin{aligned}
R_{102}&=z_{135} -z_{137}\,z_{139} -z_{138}
\end{aligned}\end{equation*}
\begin{equation*}\begin{aligned}
R_{103}&=-z_{139}\,z_{141} +z_{140} -z_{142}
\end{aligned}\end{equation*}
\begin{equation*}\begin{aligned}
R_{104}&=-z_{139}\,z_{144} +z_{143}\,z_{235} -z_{145}
\end{aligned}\end{equation*}
\begin{equation*}\begin{aligned}
R_{105}&=-z_{139}\,z_{147} +z_{146}\,z_{236} -z_{148}
\end{aligned}\end{equation*}
\begin{equation*}\begin{aligned}
R_{106}&=z_{149} -2\,z_{149}\,z_{151} +z_{152}
\end{aligned}\end{equation*}
\begin{equation*}\begin{aligned}
R_{107}&=-z_{149}\,z_{152} -z_{152} +z_{153}
\end{aligned}\end{equation*}
\begin{equation*}\begin{aligned}
R_{108}&=2\,z_{150} +z_{150}^{2} +z_{152}^{2} -z_{153}^{2}
\end{aligned}\end{equation*}
\begin{equation*}\begin{aligned}
R_{109}&=-z_{150}\,z_{155} -z_{155} +z_{156}
\end{aligned}\end{equation*}
\begin{equation*}\begin{aligned}
R_{110}&=2\,z_{154} +z_{154}^{2} +z_{155}^{2} -z_{156}^{2}
\end{aligned}\end{equation*}
\begin{equation*}\begin{aligned}
R_{111}&=z_{5}^{2} -2\,z_{5}\,z_{150} -2\,z_{150} +z_{157}
\end{aligned}\end{equation*}
\begin{equation*}\begin{aligned}
R_{112}&=-1 -z_{5}\,z_{155} +z_{7} -z_{154} \\
&\quad -z_{155} +z_{156} +z_{157}\,z_{158}
\end{aligned}\end{equation*}
\begin{equation*}\begin{aligned}
R_{113}&=1 -2\,z_{6} +z_{6}^{2} -2\,z_{6}\,z_{150} \\
&\quad +z_{159}
\end{aligned}\end{equation*}
\begin{equation*}\begin{aligned}
R_{114}&=-z_{6}\,z_{205} +z_{21} +z_{159}\,z_{160} -z_{204} \\
&\quad +z_{206}
\end{aligned}\end{equation*}
\begin{equation*}\begin{aligned}
R_{115}&=-1 +2\,z_{150}\,z_{220} -z_{161} +2\,z_{220} \\
&\quad -z_{220}^{2}
\end{aligned}\end{equation*}
\begin{equation*}\begin{aligned}
R_{116}&=-1 +z_{109} +z_{161}\,z_{162} -z_{204} \\
&\quad -z_{205}\,z_{220} +z_{206}
\end{aligned}\end{equation*}
\begin{equation*}\begin{aligned}
R_{117}&=-z_{150}\,z_{164} -z_{164} +z_{165}
\end{aligned}\end{equation*}
\begin{equation*}\begin{aligned}
R_{118}&=2\,z_{163} +z_{163}^{2} +z_{164}^{2} -z_{165}^{2}
\end{aligned}\end{equation*}
\begin{equation*}\begin{aligned}
R_{119}&=1 -2\,z_{109} +z_{109}^{2} -2\,z_{109}\,z_{150} \\
&\quad +z_{166}
\end{aligned}\end{equation*}
\begin{equation*}\begin{aligned}
R_{120}&=-1 -z_{109}\,z_{164} +z_{110} -z_{163} \\
&\quad +z_{165} +z_{166}\,z_{167}
\end{aligned}\end{equation*}
\begin{equation*}\begin{aligned}
R_{121}&=4 -4\,z_{109} +z_{109}^{2} -2\,z_{109}\,z_{150} \\
&\quad +2\,z_{150} +z_{168}
\end{aligned}\end{equation*}
\begin{equation*}\begin{aligned}
R_{122}&=-1 -z_{109}\,z_{155} +z_{111} -z_{154} \\
&\quad +z_{155} +z_{156} +z_{168}\,z_{169}
\end{aligned}\end{equation*}
\begin{equation*}\begin{aligned}
R_{123}&=1 +z_{150} -z_{150}\,z_{171} -z_{171} \\
&\quad +z_{172}
\end{aligned}\end{equation*}
\begin{equation*}\begin{aligned}
R_{124}&=1 +2\,z_{170} +z_{170}^{2} -2\,z_{171} \\
&\quad +z_{171}^{2} -z_{172}^{2}
\end{aligned}\end{equation*}
\begin{equation*}\begin{aligned}
R_{125}&=-1 +z_{109} -z_{109}\,z_{171} +z_{115} \\
&\quad +z_{166}\,z_{173} -z_{170} +z_{172}
\end{aligned}\end{equation*}
\begin{equation*}\begin{aligned}
R_{126}&=-2\,z_{150} -z_{150}^{2} +z_{174}
\end{aligned}\end{equation*}
\begin{equation*}\begin{aligned}
R_{127}&=1 -z_{109} -z_{109}\,z_{154} -z_{109}\,z_{156} \\
&\quad +z_{116} +z_{150}\,z_{156} +z_{154} -z_{155}\,z_{174} \\
&\quad +2\,z_{156} +z_{168}\,z_{175}
\end{aligned}\end{equation*}
\begin{equation*}\begin{aligned}
R_{128}&=-z_{150}\,z_{177} -z_{177} +z_{178}
\end{aligned}\end{equation*}
\begin{equation*}\begin{aligned}
R_{129}&=2\,z_{176} +z_{176}^{2} +z_{177}^{2} -z_{178}^{2}
\end{aligned}\end{equation*}
\begin{equation*}\begin{aligned}
R_{130}&=-1 -z_{109}\,z_{177} +z_{119} +z_{166}\,z_{179} \\
&\quad -z_{176} +z_{178}
\end{aligned}\end{equation*}
\begin{equation*}\begin{aligned}
R_{131}&=-z_{150}\,z_{181} -z_{181} +z_{182}
\end{aligned}\end{equation*}
\begin{equation*}\begin{aligned}
R_{132}&=2\,z_{180} +z_{180}^{2} +z_{181}^{2} -z_{182}^{2}
\end{aligned}\end{equation*}
\begin{equation*}\begin{aligned}
R_{133}&=-1 -z_{109}\,z_{181} +z_{122} +z_{166}\,z_{183} \\
&\quad -z_{180} +z_{182}
\end{aligned}\end{equation*}
\begin{equation*}\begin{aligned}
R_{134}&=-z_{150}\,z_{185} -z_{185} +z_{186}
\end{aligned}\end{equation*}
\begin{equation*}\begin{aligned}
R_{135}&=2\,z_{184} +z_{184}^{2} +z_{185}^{2} -z_{186}^{2}
\end{aligned}\end{equation*}
\begin{equation*}\begin{aligned}
R_{136}&=-1 -z_{109}\,z_{185} +z_{125} +z_{166}\,z_{187} \\
&\quad -z_{184} +z_{186}
\end{aligned}\end{equation*}
\begin{equation*}\begin{aligned}
R_{137}&=-z_{150}\,z_{189} -z_{189} +z_{190}
\end{aligned}\end{equation*}
\begin{equation*}\begin{aligned}
R_{138}&=2\,z_{188} +z_{188}^{2} +z_{189}^{2} -z_{190}^{2}
\end{aligned}\end{equation*}
\begin{equation*}\begin{aligned}
R_{139}&=-1 -z_{109}\,z_{189} +z_{128} +z_{166}\,z_{191} \\
&\quad -z_{188} +z_{190}
\end{aligned}\end{equation*}
\begin{equation*}\begin{aligned}
R_{140}&=-z_{150}\,z_{193} -z_{193} +z_{194}
\end{aligned}\end{equation*}
\begin{equation*}\begin{aligned}
R_{141}&=2\,z_{192} +z_{192}^{2} +z_{193}^{2} -z_{194}^{2}
\end{aligned}\end{equation*}
\begin{equation*}\begin{aligned}
R_{142}&=-1 -z_{109}\,z_{193} +z_{129} +z_{168}\,z_{195} \\
&\quad -z_{192} +z_{193} +z_{194}
\end{aligned}\end{equation*}
\begin{equation*}\begin{aligned}
R_{143}&=-z_{150}\,z_{197} -z_{197} +z_{198}
\end{aligned}\end{equation*}
\begin{equation*}\begin{aligned}
R_{144}&=2\,z_{196} +z_{196}^{2} +z_{197}^{2} -z_{198}^{2}
\end{aligned}\end{equation*}
\begin{equation*}\begin{aligned}
R_{145}&=-1 -z_{109}\,z_{197} +z_{132} +z_{166}\,z_{199} \\
&\quad -z_{196} +z_{198}
\end{aligned}\end{equation*}
\begin{equation*}\begin{aligned}
R_{146}&=-z_{150}\,z_{201} -z_{201} +z_{202}
\end{aligned}\end{equation*}
\begin{equation*}\begin{aligned}
R_{147}&=2\,z_{200} +z_{200}^{2} +z_{201}^{2} -z_{202}^{2}
\end{aligned}\end{equation*}
\begin{equation*}\begin{aligned}
R_{148}&=-1 -z_{109}\,z_{201} +z_{135} +z_{166}\,z_{203} \\
&\quad -z_{200} +z_{202}
\end{aligned}\end{equation*}
\begin{equation*}\begin{aligned}
R_{149}&=1 -z_{150}\,z_{205} -z_{205} +z_{206}
\end{aligned}\end{equation*}
\begin{equation*}\begin{aligned}
R_{150}&=-1 +2\,z_{204} +z_{204}^{2} +z_{205}^{2} \\
&\quad -2\,z_{206} -z_{206}^{2}
\end{aligned}\end{equation*}
\begin{equation*}\begin{aligned}
R_{151}&=-z_{109}\,z_{205} +z_{136} +z_{168}\,z_{207} -z_{204} \\
&\quad +z_{205} +z_{206}
\end{aligned}\end{equation*}
\begin{equation*}\begin{aligned}
R_{152}&=-z_{150}\,z_{209} -z_{209} +z_{210}
\end{aligned}\end{equation*}
\begin{equation*}\begin{aligned}
R_{153}&=2\,z_{208} +z_{208}^{2} +z_{209}^{2} -z_{210}^{2}
\end{aligned}\end{equation*}
\begin{equation*}\begin{aligned}
R_{154}&=-1 -z_{109}\,z_{209} +z_{140} +z_{166}\,z_{211} \\
&\quad -z_{208} +z_{210}
\end{aligned}\end{equation*}
\begin{equation*}\begin{aligned}
R_{155}&=-z_{150}\,z_{213} -z_{213} +z_{214}
\end{aligned}\end{equation*}
\begin{equation*}\begin{aligned}
R_{156}&=2\,z_{212} +z_{212}^{2} +z_{213}^{2} -z_{214}^{2}
\end{aligned}\end{equation*}
\begin{equation*}\begin{aligned}
R_{157}&=-1 -z_{109}\,z_{213} +z_{143} +z_{166}\,z_{215} \\
&\quad -z_{212} +z_{214}
\end{aligned}\end{equation*}
\begin{equation*}\begin{aligned}
R_{158}&=-z_{150}\,z_{217} -z_{217} +z_{218}
\end{aligned}\end{equation*}
\begin{equation*}\begin{aligned}
R_{159}&=2\,z_{216} +z_{216}^{2} +z_{217}^{2} -z_{218}^{2}
\end{aligned}\end{equation*}
\begin{equation*}\begin{aligned}
R_{160}&=-1 -z_{109}\,z_{217} +z_{146} +z_{166}\,z_{219} \\
&\quad -z_{216} +z_{218}
\end{aligned}\end{equation*}
\begin{equation*}\begin{aligned}
R_{161}&=-6 -2\,z_{1} -3\,z_{1}^{2} -2\,z_{1}\,z_{5} \\
&\quad -z_{2}^{2} -z_{3}^{2} -2\,z_{5} -2\,z_{5}^{2} \\
&\quad -z_{6}^{2} -z_{7}^{2} +2\,z_{12} -z_{12}^{2} \\
&\quad -z_{13}^{2} -z_{14}^{2} -z_{15}^{2} -z_{16}^{2} \\
&\quad -z_{17}^{2} -z_{18}^{2} -z_{19}^{2} -z_{20}^{2} \\
&\quad -z_{21}^{2} -z_{107}^{2} +2\,z_{109} -2\,z_{109}^{2} \\
&\quad -z_{110}^{2} -z_{111}^{2} -z_{115}^{2} -z_{116}^{2} \\
&\quad -z_{119}^{2} -z_{122}^{2} -z_{125}^{2} -z_{128}^{2} \\
&\quad -z_{129}^{2} -z_{132}^{2} -z_{135}^{2} -z_{136}^{2} \\
&\quad -z_{140}^{2} -z_{143}^{2} -z_{146}^{2} +z_{149} \\
&\quad -z_{220}^{2}
\end{aligned}\end{equation*}
\begin{equation*}\begin{aligned}
R_{162}&=-z_{6}\,z_{12} -z_{6}\,z_{13} -z_{6}\,z_{14} -z_{6}\,z_{15} \\
&\quad -z_{6}\,z_{16} -z_{6}\,z_{17} -z_{6}\,z_{18} -z_{6}\,z_{19} \\
&\quad -z_{6}\,z_{20} +z_{220}
\end{aligned}\end{equation*}
\begin{equation*}\begin{aligned}
R_{163}&=-16 +z_{1} -z_{221}^{2}
\end{aligned}\end{equation*}
\begin{equation*}\begin{aligned}
R_{164}&=-1 +z_{3} +z_{7} -z_{7}^{2}
\end{aligned}\end{equation*}
\begin{equation*}\begin{aligned}
R_{165}&=-z_{150}\,z_{223} -z_{223} +z_{231}
\end{aligned}\end{equation*}
\begin{equation*}\begin{aligned}
R_{166}&=-1 +z_{222}^{2} +z_{223}^{2} -z_{231}^{2}
\end{aligned}\end{equation*}
\begin{equation*}\begin{aligned}
R_{167}&=-z_{225} -z_{225}\,z_{226} +z_{232}
\end{aligned}\end{equation*}
\begin{equation*}\begin{aligned}
R_{168}&=-1 +z_{224}^{2} +z_{225}^{2} -z_{232}^{2}
\end{aligned}\end{equation*}
\begin{equation*}\begin{aligned}
R_{169}&=-z_{217}^{2} +z_{233}
\end{aligned}\end{equation*}
\begin{equation*}\begin{aligned}
R_{170}&=z_{223} -2\,z_{227}\,z_{233} -z_{233}
\end{aligned}\end{equation*}
\begin{equation*}\begin{aligned}
R_{171}&=z_{150}\,z_{222} +z_{228}
\end{aligned}\end{equation*}
\begin{equation*}\begin{aligned}
R_{172}&=-z_{150} -z_{222}\,z_{228} +z_{226}
\end{aligned}\end{equation*}
\begin{equation*}\begin{aligned}
R_{173}&=-1 -z_{216} -z_{222}\,z_{229} +z_{224}
\end{aligned}\end{equation*}
\begin{equation*}\begin{aligned}
R_{174}&=-z_{107} -z_{217}\,z_{230} +z_{225}
\end{aligned}\end{equation*}
\begin{equation*}\begin{aligned}
R_{175}&=2\,z_{1}^{2} +2\,z_{1}\,z_{5} -2\,z_{1}\,z_{155} -2\,z_{1}\,z_{164} \\
&\quad +z_{2}^{2} -2\,z_{2}\,z_{205} +z_{3}^{2} -2\,z_{3}\,z_{193} \\
&\quad +z_{5}^{2} -2\,z_{5}\,z_{164} +z_{12}^{2} -2\,z_{12}\,z_{171} \\
&\quad +z_{13}^{2} -2\,z_{13}\,z_{177} +z_{14}^{2} -2\,z_{14}\,z_{181} \\
&\quad +z_{15}^{2} -2\,z_{15}\,z_{185} +z_{16}^{2} -2\,z_{16}\,z_{189} \\
&\quad +z_{17}^{2} -2\,z_{17}\,z_{197} +z_{18}^{2} -2\,z_{18}\,z_{201} \\
&\quad +z_{19}^{2} -2\,z_{19}\,z_{209} +z_{20}^{2} -2\,z_{20}\,z_{213} \\
&\quad +z_{107}^{2} -2\,z_{107}\,z_{217} -z_{150}\,z_{234} -z_{154} \\
&\quad +z_{155}^{2} -z_{163} +z_{164}^{2} -z_{170} \\
&\quad +z_{171}^{2} -z_{176} +z_{177}^{2} -z_{180} \\
&\quad +z_{181}^{2} -z_{184} +z_{185}^{2} -z_{188} \\
&\quad +z_{189}^{2} -z_{192} +z_{193}^{2} -z_{196} \\
&\quad +z_{197}^{2} -z_{200} +z_{201}^{2} -z_{204} \\
&\quad +z_{205}^{2} -z_{208} +z_{209}^{2} -z_{212} \\
&\quad +z_{213}^{2} -z_{216} +z_{217}^{2}
\end{aligned}\end{equation*}
\begin{equation*}\begin{aligned}
R_{176}&=-z_{111}\,z_{235} +z_{136}
\end{aligned}\end{equation*}
\begin{equation*}\begin{aligned}
R_{177}&=-z_{116}\,z_{236} +z_{136}
\end{aligned}\end{equation*}
\begin{equation*}\begin{aligned}
R_{178}&=-z_{129}\,z_{237} +z_{136}
\end{aligned}\end{equation*}
\begin{equation*}\begin{aligned}
R_{179}&=-2 -z_{107}^{2} -z_{155}^{2} -z_{164}^{2} \\
&\quad +2\,z_{171} -z_{171}^{2} -z_{177}^{2} -z_{181}^{2} \\
&\quad -z_{185}^{2} -z_{189}^{2} -z_{193}^{2} -z_{197}^{2} \\
&\quad -z_{201}^{2} -z_{205}^{2} -z_{209}^{2} -z_{213}^{2} \\
&\quad +z_{217}^{2} -z_{238}^{2} -z_{239}^{2} -z_{240} \\
&\quad -z_{240}^{2}
\end{aligned}\end{equation*}
\begin{equation*}\begin{aligned}
R_{180}&=-1 -4\,z_{113}^{2} -4\,z_{118}^{2} -4\,z_{121}^{2} \\
&\quad -4\,z_{124}^{2} -4\,z_{127}^{2} -4\,z_{131}^{2} -4\,z_{134}^{2} \\
&\quad +z_{136}^{2} -4\,z_{138}^{2} -4\,z_{142}^{2} -4\,z_{145}^{2} \\
&\quad -4\,z_{148}^{2} -z_{241}^{2} -z_{242}^{2} -z_{243} \\
&\quad -z_{243}^{2}
\end{aligned}\end{equation*}
\begin{equation*}\begin{aligned}
R_{181}&=-10 +z_{109} -36\,z_{244} -36\,z_{244}^{2}
\end{aligned}\end{equation*}
\begin{equation*}\begin{aligned}
R_{182}&=z_{6}\,z_{8} -z_{8} -z_{8}^{2} -z_{245}^{2} \\
&\quad -z_{246}^{2} -z_{247} -z_{247}^{2}
\end{aligned}\end{equation*}

\endgroup

\clearpage
\section{Exact data format and integrity}

In a polynomial JSON file a term is a pair
$[c,[i_1,\ldots,i_d]]$, denoting
$c\,z_{i_1+1}\cdots z_{i_d+1}$. Indices are zero-based in the files
and one-based in the displayed mathematics. An empty index list is
a constant, repeated indices are powers, all coefficients are integers,
and all like monomials have been collected in the expansion.

The SHA-256 digest of the final residual equation file is
\begin{quote}\small\ttfamily
6df8f7de8b389084d6ce000f777f27ffc962e05b51217662e680201cffd1578d
\end{quote}
and that of the collected polynomial is
\begin{quote}\small\ttfamily
19b2bc4e8cbb8d112feed0b88ca6c91884a40985745026eed020a78eee35187f.
\end{quote}
The period after the second digest is punctuation. The raw
247-variable weak-product equation before the final integral maps has
digest
\begin{quote}\small\ttfamily
14b814894521add09fb6847e59de2d802f2e466f6f9fc724f346a21461125a6b.
\end{quote}
The hashes bind finite files; they do not constitute a mathematical
proof. The manifest separately binds all deliverable files, excluding
the manifest itself.
\end{document}
